% !TeX root = main.tex
\begin{abstract}
대규모 언어 모델(LLM)은 협력형 다중 에이전트 강화학습에서 장기 목표를 서브골과
행동 제약으로 분해하는 고수준 의사결정기로 활용될 수 있다. 대표적으로 LEHCA는
LLM Commander가 생성한 의미적 지침을 환경 전이에서 계산 가능한 보조 신호와
행동 규칙으로 변환하고, 저수준 에이전트가 이를 이용해 정책을 학습하도록 한다.
이처럼 언어 서브골을 실행 가능한 신호에 대응시키는 과정을 의미 접지(semantic
grounding)라 한다. 그러나 LLM 지침은
저수준 행동보다 느린 시간척도에서 갱신되므로, 언제 새 지침을 요청할 것인지가 별도의
의사결정 문제가 된다. 짧은 고정 주기는 불필요한 호출과 출력 변동을 늘리고, 긴 고정
주기는 과업 상태나 달성 가능한 서브골이 변한 뒤에도 오래된 지침을 유지할 수 있다.

본 연구는 고정된 호출 간격 대신 현재 보유 지침이 앞으로 생성할 것으로 예상되는
셰이핑 신호를 이용해 갱신 시점을 결정하는 RSVP(Residual Shaping-Value Prediction)를
제안한다. RSVP는 각 접지 함수가 생성할 보조 신호의 단기 할인 누적값을 다중 출력
모니터링 모델로 예측한다. 현재 지침의 가중치로 합성한 잔여 셰이핑 가치가 발급 시점에
비해 지속적으로 감소하면 단측 CUSUM 통계량을 통해 조기 갱신한다. 최대 보유 간격은
안전한 강제 갱신 기준으로 유지하므로, 검출이 불가능한 상황에서도 지침이 무기한
고착되지 않는다. 이 방법은
별도의 LLM 평가 호출 없이 기존 셰이핑 경로에서 계산되는 신호를 재사용한다.
현재 연구의 핵심 검증 과제는 기반 LEHCA 지침 채널의 유효성, 잔여 가치의 시간순
예측력, 그리고 검출 시점에서 실제 갱신이 외부 성능을 향상시키는지를 각각 분리해
확인하는 것이다.
\end{abstract}
